% ============================================================
% CAPITOLO 28 - SETTIME DI QUARTA E QUINTA SPECIE
% ============================================================
\chapter{Settime di quarta e quinta specie}
\label{chap:quarta_quinta_specie}

\section{Quarta specie: uso e risoluzione}
\label{sec:quarta_specie_dettaglio}

La quarta specie è formata da una \textbf{triade maggiore} e una \textbf{settima maggiore}. Nel corso di armonia si studia solo l'accordo costruito sul \textbf{primo grado} del modo maggiore.

L'accordo ha una tensione particolare e risolve di norma sulla \textbf{sottodominante (IV)}. La settima, che è la sensibile della scala, è preparata e scende di grado verso la terza del IV grado.

\begin{esempio}[\sigla{I}{7}{} della quarta specie]
    \musica{es_cap13_maggiore_do_1}
    \musica{es_cap13_maggiore_do_2}
    Il \sigla{I}{7}{} risolve sul IV: la settima Si scende a La. Il primo rivolto (basso Mi) sale a Fa; il secondo (basso Sol) scende a Fa; il terzo (basso Si, la settima) scende a La sul \sigla{IV}{6}{}.
\end{esempio}

\begin{esempio}[Quarta specie in La minore: \sigla{VI}{7}{} $\rightarrow$ II]
    \label{ex:vi7_minore}
    Sul VI grado della minore armonica la settima maggiore (Fa-La-Do-Mi) risolve per quinta discendente sul II grado (triade diminuita Si-Re-Fa): la settima Mi scende a Re.
    \musica{es_cap13_maggiore_la_1}
    \musica{es_cap13_maggiore_la_2}
\end{esempio}

\section{Quarta specie: collocazione e rivolti}
\label{subsec:iv7_collocazione}

\begin{regola}[Gradi della quarta specie]
	\label{reg:iv7_gradi}
	L'accordo di settima maggiore (quarta specie) si trova:
	\begin{itemize}
		\item sul \textbf{I} e sul \textbf{IV} grado della scala maggiore;
		\item sul \textbf{VI} grado della scala minore armonica;
		\item sul \textbf{III} e sul \textbf{VI} grado della scala minore melodica discendente.
	\end{itemize}
	L'analisi del corso si concentra sul \textbf{I grado del modo maggiore}.
\end{regola}

Sul I grado l'accordo di quarta specie conserva la funzione di tonica, ma la settima lo spinge verso il \textbf{IV grado}, su cui risolve (sul IV grado lo stesso tipo di accordo ha funzione di sottodominante). La settima va preparata nella voce; le altre note sono parti libere, e la quinta si può omettere a favore del raddoppio della fondamentale.

\begin{center}
\begin{tabularx}{\linewidth}{lLLLL}
\toprule
\textbf{Stato} & \textbf{Cifratura} & \textbf{Grado al basso} & \textbf{Risolve su} & \textbf{Basso obbligato} \\
\midrule
Fond. & 7 & I & IV (o \sigla{IV}{7}{}) & no \\
I rv. & 6/5 & III & IV & sì (sale) \\
II rv. & 4/3 & V & IV & sì (scende) \\
III rv. & 4/2 & VII (settima) & \sigla{IV}{6}{} & sì (scende) \\
\bottomrule
\end{tabularx}
\end{center}

\section{Quinta specie: natura e collocazione}
\label{sec:quinta_specie_dettaglio}

La quinta specie è la \textbf{settima diminuita} (triade diminuita e settima diminuita). È l'accordo più dissonante e più simmetrico.

\begin{regola}[Grado della quinta specie]
	\label{reg:vii_dim7_gradi}
	La settima diminuita si trova sul \textbf{VII grado della scala minore armonica} (nella melodica ascendente il VII grado dà invece una semidiminuita), dove si chiama \textbf{settima di sensibile del modo minore}. È l'unico accordo di settima che contiene l'intervallo di \textbf{settima diminuita}, e si riconosce bene all'ascolto.
\end{regola}

È un accordo \textbf{derivato}: nasce dall'accordo di nona di dominante minore senza la fondamentale e, per questo, \textbf{non richiede la preparazione} della settima.

\section{Uso sul VII grado del minore}
\label{subsec:vii_dim7_natura}

Si studia soprattutto sul \textbf{VII grado del modo minore} (sulla sensibile), dove ha funzione di dominante estrema. Alla risoluzione sul I grado valgono le norme della settima di sensibile di terza specie (Capitolo \ref{chap:vii7_ii7}): la sensibile sale alla tonica, la quinta e la settima scendono per grado.

\begin{esempio}[\sigla{VII°}{7}{} $\rightarrow$ I in La minore]
	\label{ex:vii_dim7}
	La sensibile (Sol$\sharp$) sale alla tonica, la quinta diminuita (Re) e la settima (Fa) scendono per grado.
	\musica{es_cap13_dim_la_1}
	\musica{es_cap13_dim_la_2}
	Nei rivolti il basso segue il suo movimento obbligato: nel \sigla{VII}{6}{5} (basso Si) scende a La, nel \sigla{VII}{4}{3} (basso Re) scende a Do, nel \sigla{VII}{2}{} (basso Fa, la settima) scende a Mi sul \sigla{i}{6}{4}. La sensibile sale sempre al La.
\end{esempio}

I rivolti ricalcano quelli della settima di sensibile (terza specie sul VII), con una sola eccezione: il terzo rivolto.

\begin{osservazione}[Terzo rivolto senza preparazione]
	\label{oss:dim7_terzo_rivolto}
	A differenza della terza specie, nella quinta specie il \textbf{terzo rivolto} si può disporre \textbf{senza preparare la settima al basso}. Il motivo teorico è che l'urto di \textbf{seconda aumentata} tra la fondamentale (sensibile) e la settima è \textbf{enarmonico} con la terza minore, consonante, e all'orecchio risulta privo della durezza delle seconde.
\end{osservazione}

Oltre che sul VII grado, la quinta specie si usa come \textbf{accordo alterato sul IV grado} di entrambi i modi, come «dominante della dominante» o per colorire le cadenze.
